\documentclass[10pt]{article}
\usepackage[margin=20mm]{geometry}
\usepackage{amsmath,amssymb,mathtools,fontspec,unicode-math,xurl,hyperref}
\IfFontExistsTF{DejaVu Serif}{\setmainfont{DejaVu Serif}}{\setmainfont{Latin Modern Roman}}
\setmathfont{latinmodern-math.otf}
\setlength{\parindent}{0pt}\setlength{\parskip}{6pt}\sloppy
\title{Scientific review of
Stability for the Sobolev inequality:
Existence of a minimizer}
\author{AI-assisted technical review}
\date{26 October 2026}
\begin{document}
\maketitle
Abstract: The assigned arXiv v4 is examined through the two-bubble expansion and the compactness argument. The literal compactness statement of Theorem 1.2 is false without translations and dilations: a normalized minimizing sequence can be translated to have weak limit zero. This is a Major statement correction, even though compactness modulo these symmetries and the intended existence conclusion remain supported by a repair of the proof. A second Major issue concerns invalid limiting divisions in the balancing argument on the route to compactness; the explicit repair does not make that core proof gap peripheral. Four Minor groups concern the zero-function boundary and approximate identity, a dimension-dependent remainder, derivative constants, and an asymptotic spectral bound stated pointwise. Decisive examples and explicit repairs are given. The assessment is Major revision for this precise author version, with scientific adjudication and human review pending.\par
\section*{Identity, definitions and theorem scope}
The author is Tobias König. The journal association is DOI https://doi.org/10.4171/JEMS/1582. The actual reviewed input is the 18-page author preprint arXiv:2211.14185v4, uploaded on 6 October 2023 and dated 5 October 2023 in the manuscript. The page numbers below are both PDF and printed page numbers. EMS records submission on 22 December 2022, acceptance on 30 September 2023, and online-first publication on 3 January 2025. The publisher PDF is subscription-restricted. No equivalence of the reviewed author file with the Version of Record has been established.\par
Put \[p=\frac{2d}{d-2s}>2,\qquad a=\frac{d-2s}{2}>0,\qquad
 H(f)=\|(-\Delta)^{s/2}f\|_2^2.\] The functional is \[\mathcal E(f)=\frac{H(f)-S_d\|f\|_p^2}{\operatorname{dist}(f,\mathcal M)^2},
 \qquad
 \mathcal M=\{c(a_0+|x-b_0|^2)^{-(d-2s)/2}:a_0>0,\ b_0\in\mathbb R^d,\ c\ne0\}.\] The different symbols \(a\) and \(a_0\) distinguish the exponent from the bubble parameter. The setting is real-valued functions; arbitrary complex phases or positivity of an eventual minimizer are not additional conclusions of the paper.\par
Theorem 1.1 asserts \(c_{\rm BE}(s)<2-2^{2/p}\) for every \(d\ge1\), \(0<s<d/2\). The existence argument uses in addition the strict spectral comparison \(c_{\rm BE}(s)<4s/(d+2s+2)\), cited from the author’s earlier paper, whose range is \(d\ge2\). Consequently no existence theorem in dimension one, endpoint statement at \(s=0,d/2\), explicit value of the best constant, or uniqueness of a minimizer follows here.\par
\section*{Established problems and repair boundaries}
Problem 1 (Major: Theorem 1.2 omits the noncompact symmetries).  Location. Theorem 1.2 and the preceding assertion on p. 4; the proof’s change of sequence at the bottom of p. 10 and top of p. 11.\par
The theorem states that every minimizing sequence with \(H(u_n)=1\) has a strongly convergent subsequence in \(\dot H^s(\mathbb R^d)\). Its statement includes neither translations and dilations nor a centering condition. The proof, however, explicitly translates and rescales the sequence before extracting a nonzero weak limit. These are changes to the sequence, not subsequence extraction.\par
A counterexample to the literal statement does not require assuming that a minimizer already exists. Take any normalized minimizing sequence \(f_n\). Let \(\{\phi_j\}_{j\ge1}\) be a countable dense set in \(\dot H^s(\mathbb R^d)\). For each fixed \(n,j\), \[\langle f_n(\,\cdot-y),\phi_j\rangle_{\dot H^s(\mathbb R^d)}\longrightarrow0
 \quad\text{as }|y|\longrightarrow\infty.\] Indeed, the Fourier integrand \(\lvert\xi\rvert^{2s}\widehat f_n(\xi)\overline{\widehat\phi_j(\xi)}\) is integrable by Cauchy–Schwarz, so the Riemann–Lebesgue lemma applies. Choose \(y_n\) so that \[|\langle f_n(\,\cdot-y_n),\phi_j\rangle_{\dot H^s(\mathbb R^d)}|\le n^{-1}
 \quad (j\le n),
 \qquad v_n=f_n(\,\cdot-y_n).\] Translation invariance gives \(H(v_n)=1\) and \(\mathcal E(v_n)=\mathcal E(f_n)\). Density of the test set and uniform boundedness imply \(v_n\rightharpoonup0\). Every subsequence has that same weak limit, whereas a strongly convergent subsequence would have a limit of norm one. Thus no such subsequence exists.\par
Impact and repair boundary. This contradicts a quantified conclusion of a main theorem. The required statement is that, after a subsequence, there are \(\lambda_n>0\) and \(x_n\in\mathbb R^d\) such that \[\lambda_n^{(d-2s)/2}u_n(\lambda_n\,\cdot+x_n)\longrightarrow u
 \quad\text{strongly in }\dot H^s(\mathbb R^d),\qquad H(u)=1,\quad \mathcal E(u)=c_{\rm BE}(s).\] The proof on pp. 10–15 is directed toward this corrected conclusion. Its use of \(L^p\) normalization can be converted back to energy normalization because the normalization factors have a positive limiting value. The existence conclusion is preserved by this repair, subject to the local proof corrections below and the cited compactness inputs. The Major classification concerns the false unqualified compactness assertion; it does not assert that the existence theorem is false or irreparable.\par
Problem 2 (Minor: the zero-function boundary and the approximate-identity proof).  Location. Definitions (1.2), (1.4)–(1.5), pp. 1,3; Lemma 2.2 and its last paragraph, pp. 6–7.\par
Because \(c=0\) is excluded from \(\mathcal M\), the stated domain \(\dot H^s(\mathbb R^d)\setminus\mathcal M\) contains zero. At zero, the functional is \(0/0\). Also, \[\inf_{h\in\mathcal M}H(h)=0\] is not attained, contradicting the unconditional attainment assertion in Lemma 2.2. The supremum \(m(0)=0\) is attained by every normalized bubble; its corresponding coefficient is zero, and the resulting function is not in the stated \(\mathcal M\). The square-completion identity itself remains true as an identity of infima.\par
There is a separate error at the top of p. 7: the approximate-identity asymptotic has \(\lambda^{+(d-2s)/2}\), although p. 6 has the correct negative exponent. In fact, with \(K=B^{p-1}\in L^1(\mathbb R^d)\) and \(K_\lambda(x)=\lambda^dK(\lambda x)\), \[(f,B_{x,\lambda}^{p-1})=\lambda^{-a}(K_\lambda*f)(x).\] A positive power is already incompatible with the uniform Hölder bound \(\lvert(f,B_{x,\lambda}^{p-1})\rvert\le\|f\|_p\) at a point where the asserted leading coefficient is nonzero. The passage from continuous approximants to a pointwise asymptotic also needs an approximate-identity argument, rather than an unqualified exchange of the approximation and concentration limits.\par
Impact and repair boundary. Define \(\mathcal E\) on \(\dot H^s(\mathbb R^d)\setminus(\mathcal M\cup\{0\})\), and restrict distance attainment to \(f\ne0\), or state an explicit convention using the closed cone \(\mathcal M\cup\{0\}\). For the needed positivity, no pointwise exchange is necessary: if \(m(f)=0\), then \(K_\lambda*f=0\) for every \(\lambda\); convergence in \(L^p\) to \((\int K)f\) implies \(f=0\). This proves \(m(f)>0\) for nonzero \(f\), and escape of bubble parameters then cannot maximize it. These corrections preserve the distance formula and all its uses for nonzero functions.\par
Problem 3 (Major: missing boundary cases in the balancing argument).  Location. Lemma 2.4, p. 7; the definitions of \(A,B,C,D\) and the limiting argument in Lemma 4.3, pp. 13–14; the quotient involving \(f\) before (4.14), p. 14.\par
The second strictness clause in Lemma 2.4 says that \(D<F\) alone makes \((A+C)/(B+D)>(A+E)/(B+F)\). Set \[A=B=C=D=1,\qquad E=F=2.\] Every hypothesis holds, including \(D<F\), but both displayed quotients and all three initial ratios equal one. This exact rational counterexample was also checked in an original calculation. The correct strictness condition is \[C/D>E/F
 \quad\text{or}\quad
 \bigl(D<F\ \text{and}\ A/B>C/D\bigr).\] The application on p. 13 does obtain the first condition, so this error alone does not invalidate that application.\par
The same passage needs two further corrections. First, \(B\) is printed as \(\lim[H(\widetilde g_n)-S_dm(\widetilde g_n)^2]\), whereas (2.3) and the immediately preceding denominator require \[B=\lim[H(\widetilde g_n)-S_dm(\widetilde g_n)].\] The quantity \(m\) is already a squared pairing. Second, being outside \(\mathcal M\) at every index does not make the limiting deficit and distance strictly positive: normalized perturbations of a bubble can have both limits zero. Thus \(A/B\) may be undefined. The separate treatment of a sequence lying exactly in \(\mathcal M\) does not address nonbubble sequences approaching it. In the reversed case, positivity of the energy alone does not justify a bounded rescaling to match the two \(m\)-values when the smaller one tends to zero. On p. 14 the weak limit \(f\) may itself be a bubble, again making the displayed quotient \(0/0\).\par
Severity and actual dependency. The unused strictness clause and the extra square in \(B\) are local errors. The limiting divisions are different: Lemma 4.3 supplies the balance of the two components used immediately in (4.14), and that estimate excludes residual mass in the proof of strong compactness and existence. Invoking strong convergence to exclude a bubble weak limit at this point would be circular. The written proof therefore has a Major gap in a core derivation. The examples diagnose the invalid inference, not a counterexample satisfying all hypotheses of the full minimizing-sequence argument. The repair below supports the intended conclusion without downgrading the original gap.\par
Impact and repair boundary. Work throughout with the continuous inequality \[H(v)-S_d\|v\|_p^2\ge c_{\rm BE}(s)\operatorname{dist}(v,\mathcal M)^2,\] which remains meaningful on a bubble. Write \(c=c_{\rm BE}(s)\). In the limiting notation of p. 13 it yields \(A\ge cB\) with \(A,B\ge0\). Since \(c(B+D)=A+C\), one gets \(C\le cD\). If the smaller component is enlarged so that \(F>D\) and \(E/F<C/D\), direct subtraction gives \[\begin{aligned}
 c(B+F)-(A+E)
 &=c(F-D)+C-E\\
 &>(F-D)(c-C/D)\ge0.
 \end{aligned}                                                    \tag{2.1}\] This is the desired strict decrease without dividing by \(A\) or \(B\).\par
The enlargement can be a fixed factor \(t>1\), sufficiently near one that the strict ordering of the two \(m\)-values is preserved. It need not continue all the way to equality. Strict increase of \[h(\eta)/\eta^2,\qquad h(\eta)=(1+\eta^p)^{2/p}-1,\] then supplies \(E/F<C/D\). This also avoids an unbounded scaling factor if the smaller limiting \(m\)-value is zero.\par
For completeness, in the reversed case the remainder cannot have vanishing \(L^p\) norm while retaining positive energy. If \(\|g_n\|_p\to0\), then \(m(g_n)\to0\). The splitting identities and the continuous stability inequality for \(f\) give \[\bigl[H(u_n)-S_d\|u_n\|_p^2\bigr]
 -c\,\operatorname{dist}(u_n,\mathcal M)^2
 \ge (1-c)H(g_n)+o(1).\] The left side tends to zero for a minimizing sequence, whereas the right side stays positive when \(H(g_n)\ge\delta>0\), because \(c<1\). Thus the remaining \(L^p\) mass is positive, and the fixed-factor comparison is strict. The same continuous inequality treats \(f\in\mathcal M\) when deriving (4.14). These steps repair the written limit argument and support its conclusion; they are not counterexamples to Lemma 4.3 or to existence.\par
Problem 4 (Minor: the quadratic remainder estimate misses two parameter regimes).  Location. Proof of Proposition 3.1(ii), p. 9, the display estimating \(\int_{B(0,\lambda^{-1/2})}B_\lambda^2B^{p-2}\).\par
The claimed \(O(\lambda^{d/2})\) estimate is false when \(d\le4s\). On this ball, for small \(\lambda\), \[I_\lambda:=\int B_\lambda^2B^{p-2}
 \asymp \lambda^{d-2s}
       \int_0^{\lambda^{-1/2}}r^{d-1}(1+r^2)^{-2s}\,dr.\] Consequently the correct orders are \[I_\lambda\asymp
 \begin{cases}
 \lambda^{d/2},&d>4s,\\
 \lambda^{d/2}\log(1/\lambda),&d=4s,\\
 \lambda^{d-2s},&d<4s.
 \end{cases}                                                       \tag{2.2}\] For a concrete allowed case, \(d=3,s=1\), normalization gives \(c_d^6=4/\pi^2\), and dominated convergence in the radial integral yields \[\lim_{\lambda\downarrow0}\frac{I_\lambda}{\lambda}
 =c_d^6\,4\pi\int_0^\infty\frac{r^2}{(1+r^2)^2}\,dr=4.\] Thus \(I_\lambda\) cannot be \(O(\lambda^{3/2})\).\par
Impact and repair boundary. Each line of (2.2) is still \(o(\lambda^{(d-2s)/2})\) for fixed \(0<s<d/2\). Replace the overly strong bound by the appropriate case, or prove just this common little-oh estimate. The first-order two-bubble expansion and its strictly negative correction remain intact over the stated parameter range. No uniform assertion as \(s\) approaches an endpoint is required.\par
Problem 5 (Minor: the exact derivative constants in Appendix A are incorrect).  Location. Lemma A.1, formula (A.1) and its derivation, pp. 15–16; use in (3.5), p. 10.\par
Retain \(a=(d-2s)/2\), and put \(b=(d+2s)/2\) and \(\omega_{d-1}=|\mathbb S^{d-1}|\). From definition (3.3), the actual normalized correlation is \[F(\mu)=c_d^p\omega_{d-1}\mu^b
 \int_0^\infty r^{d-1}(1+r^2)^{-a}(1+\mu^2r^2)^{-b}\,dr.\] Differentiating this expression, or using \(r^2/(1+r^2)\) with beta distribution of parameters \(d/2,d/2\), gives \[F'(1)=0,\qquad
 F''(1)=-\frac{ab}{d+1}
       =-\frac{d^2-4s^2}{4(d+1)}.\] Also, direct differentiation and dominated convergence, uniformly for \(\mu\) near one, give \[G_\lambda(\mu)=c_0\lambda^a\mu^{-a}+o(\lambda^a),
 \qquad G_\lambda'(1)=-a c_0\lambda^a+o(\lambda^a).\] Both constants printed in (A.1) are missing the common factor \(\omega_{d-1}b\). The intermediate differentiation formulas need correction as well: the radial surface measure and the derivative of \(\mu^b\) must be retained, and the coefficient inside the resulting bracket is \(1/a=2/(d-2s)\), rather than \(2/(d+2s)\).\par
An exact elementary case makes the discrepancy explicit. For \(d=3,s=1/2\), one has \(p=3\), \(c_d^3=4/\pi^2\), and \[F(\mu)=\frac{4\mu}{(1+\mu)^2},\qquad
 G_\lambda(\mu)=\frac{4\lambda\mu}{(\lambda+\mu)^2}.\] Hence \(F''(1)=-1/2\), \(c_0=4\), and \(G_\lambda'(1)=-4\lambda+o(\lambda)\). The printed constants are instead \(-1/(16\pi)\) and \(-1/(2\pi)\). The elementary radial integral and exact rational derivatives were independently checked; no author code was executed.\par
Impact and repair boundary. The actual constants still have the needed sign and size: \(F''(1)<0\) and \(G_\lambda'(1)=O(\lambda^a)\). Re-differentiating the correct normalized integral repairs the local maximum argument, including \(\mu(\lambda)-1=O(\lambda^a)\), and preserves Proposition 3.1(iii). The error concerns precise formulas, not an established failure of the two-bubble theorem.\par
Problem 6 (Minor: the spectral lower bound is asymptotic, not pointwise).  Location. Proof of Proposition 4.1, p. 11, the assertion \(\mathcal E(u_n)\ge4s/(d+2s+2)\) for sequences approaching \(\mathcal M\).\par
For the normalized sequences in question the needed statement is \[\liminf_{n\to\infty}\mathcal E(u_n)\ge\kappa,
 \qquad \kappa=\frac{4s}{d+2s+2}.\] A pointwise inequality need not hold. The already cited earlier paper [sharp], Theorem 1 and pp. 3–5 of its arXiv v2, constructs perturbations \(U+\varepsilon\rho\) approaching a bubble with quotient strictly below \(\kappa\) for every sufficiently small positive \(\varepsilon\).\par
A concrete direction in the classical case \(d=3,s=1\) comes from the degree-two spherical harmonic \[v_2(\omega)=\omega_1\omega_2+\omega_2\omega_3+\omega_3\omega_1
 \quad\text{on }\mathbb S^3.\] Its cubic moment is \(6\int_{\mathbb S^3}\omega_1^2\omega_2^2\omega_3^2>0\). The corresponding normal perturbation has quotient \(\kappa-K\varepsilon+o(\varepsilon)\), \(K>0\). Dividing by its \(L^p\) norm preserves the quotient and gives the normalization used in Section 4, while the distance to \(\mathcal M\) tends to zero. Thus it also demonstrates the distinction within the relevant normalization.\par
Impact and repair boundary. Replace the pointwise assertion by the liminf bound. The strict comparison \(c_{\rm BE}(s)<\kappa\) then gives exactly the contradiction needed in Proposition 4.1. The underlying spectral and nearest-bubble inputs remain the stated external dependencies; a full independent reconstruction of those foundational results is not claimed.\par
\section*{Positive checks and proof dependencies}
\subsection*{The two-bubble threshold.}
The normalized bubble satisfies \(B^{p-1}\in L^1\), since its decay exponent is \(d+2s>d\); hence \(c_0>0\) is finite. The pairing identity from the Euler–Lagrange equation gives \[H(B+B_\lambda)=2S_d+2S_dc_0\lambda^a+o(\lambda^a).\] The inversion of radius \(\lambda^{-1/2}\) exchanges the two bubbles. Splitting the \(L^p\) integral at that radius, using the corrected remainder (2.2), yields \[\|B+B_\lambda\|_p^2
 =2^{2/p}+2^{2/p+1}c_0\lambda^a+o(\lambda^a).\] Positivity and radial decrease justify centering the maximizing comparison bubble at zero. The symmetry \(H_\lambda(\mu)=H_\lambda(\lambda/\mu)\) reduces scale optimization to \(\mu\ge\sqrt\lambda\). The estimates (3.4) keep maximizing scales in a compact interval away from zero; Hölder equality forces their limit to be one. Correct differentiation then gives \[m(B+B_\lambda)^{1/2}=1+c_0\lambda^a+o(\lambda^a).\] The first-order term in the distance cancels, so \[\operatorname{dist}(B+B_\lambda,\mathcal M)^2=S_d+o(\lambda^a),\qquad
 \mathcal E(B+B_\lambda)
 =2-2^{2/p}-(2^{2/p+1}-2)c_0\lambda^a+o(\lambda^a).\] The coefficient subtracted is strictly positive for every allowed \(p\). This verifies the mechanism of Theorem 1.1 after the stated local corrections, including dimension one.\par
\subsection*{Splitting and exclusion of residual mass.}
The energy identity follows from weak convergence in a Hilbert space. Local compactness followed by Brezis–Lieb gives the \(L^p\) mass identity after a subsequence. The maximum, rather than a sum, in Lemma 4.2 is essential and is supported by its two cases. A comparison bubble whose parameters remain compact pairs negligibly with \(g_n\rightharpoonup0\); a bubble whose parameters escape pairs negligibly with the fixed \(f\). Testing separately with optimizers for \(m(f)\) and \(m(g_n)\) gives the reverse inequality. This establishes \[m(f+g_n)=\max\{m(f),m(g_n)\}+o(1).\] The squared-pairing definition accommodates either sign of a component; no positivity assumption on a general minimizing sequence is needed.\par
With the corrected comparison in Problem 3, a nonvanishing remainder forces \(m(f)=m(g_n)+o(1)\). The larger \(L^p\) component may then be normalized to one, and the smaller has norm at most \(1+o(1)\). The continuous stability inequality for the larger component and strict increase of \(h(\eta)/\eta^2\) give the limiting form of (4.14), \[c_{\rm BE}(s)\ge
 1-\frac{S_d}{S(g_n)}(2^{2/p}-1)+o(1)
 \ge 2-2^{2/p}+o(1),\] because \(S(g_n)\ge S_d\). The positive gap in Theorem 1.1 excludes this possibility. A strong limit lies outside \(\mathcal M\) by the asymptotic spectral bound and its strict global comparison. These arguments support existence and compactness modulo translations and dilations, not the unqualified statement refuted in Problem 1.\par
The initial extraction of a nonzero profile uses the cited critical compactness theorem. The global fractional stability inequality, its local spectral expansion, and the compactness theorem are used as external inputs. The earlier strict spectral construction was read in its six-page arXiv v2 as a bounded dependency and counterevidence check; its cubic direction and dimension restriction were examined. Reading those six pages is not an independent proof of all their cited spectral or geometric results.\par
\section*{Coverage, attribution and assessment}
During the initial review, the principal read all 941 lines of the supplied primary text extraction and actually inspected all 18 primary PDF page images. It also read all 305 lines and inspected all six pages of the separately retrieved arXiv:2210.08482v2 used for the bounded check of reference [21]. The latter is an author version of a different, cited paper; it is not a substitute for the assigned v4 or for either publisher version.\par
Their actual page and TeX-line coverage, input hashes and process exits are retained in their complete memoranda. The principal evaluated the findings, wrote this English TeX report, executed the existing compiler, and inspected every final report page. The auxiliaries are not independent acceptance reviewers or independently authenticated Pro backends.\par
The exact rational checks concern the scalar counterexample and the closed-form \(d=3,s=1/2\) correlation. The radial integral and functional-analytic counterexample are mathematical derivations; they are not represented as numerical experiments. No author code or author TeX was executed. A bounded publisher, arXiv and correction/supplement search found no separate correction or executable supplement for this exact input. Lack of code is not a mathematical defect.\par
There are six problem groups: two Major issues, concerning the main compactness statement and the core balancing derivation, and four Minor groups with explicit repair boundaries. Focused independent acceptance rechecked each reported group, its source location, decisive evidence, counterevidence, actual proof dependency, and repair scope. The present revision corrects the severity and explanation of Problem 3; native assistance was limited to read-only challenges and did not constitute human review or a separately authenticated backend. Substantive support remains for the intended existence mechanism after the stated repairs. No conclusion about an uninspected publisher version is drawn. The assessment is Major revision for the precise assigned author version. Scientific adjudication remains `scientific\_decision=pending`, with `human\_reviews=[]`.\par
9 T. König, Stability for the Sobolev inequality: existence of a minimizer, arXiv:2211.14185v4, uploaded 6 October 2023, manuscript dated 5 October 2023. https://arxiv.org/abs/2211.14185v4. EMS Press, article record, DOI 10.4171/JEMS/1582. https://ems.press/journals/jems/articles/14298462. T. König, On the sharp constant in the Bianchi–Egnell stability inequality, Bull. London Math. Soc. 55 (2023), 2070–2075, DOI https://doi.org/10.1112/blms.12837. Bounded source checked: arXiv:2210.08482v2, 6 November 2022. https://arxiv.org/pdf/2210.08482v2.\par
\end{document}
